%% Chapter 7 -----------------------------------------------------------------
\chapter{Viewing, Processing and Annotation Tools}
\label{ch:processing}

This chapter describes how to navigate the spectrum, set an exact display range, clean
radio-frequency interference, use processing presets, annotate the plot, overlay light
curves, and make quick drift-rate and ruler measurements.

\section{Navigating the spectrum}
\label{sec:navigation}
\index{navigation}\index{zoom}

Table~\ref{tab:navigation} lists the mouse interactions available in the plotting area.

\begin{table}[!htbp]
  \centering
  \caption{Navigation and picking in the plotting area.}
  \label{tab:navigation}
  \small
  \begin{tabularx}{\linewidth}{@{}P{0.27\linewidth}L@{}}
    \toprule
    \thead{Interaction} & \thead{How} \\
    \midrule
    Zoom & Scroll the mouse wheel over the plot; the zoom is centered on the pointer. \\
    Pan & Hold the left mouse button and drag. \\
    Rectangle zoom & Click \gui{Lock zooming and panning}\index{navigation!lock and unlock} on the toolbar, then \gui{Rectangular Zooming}\index{zoom!rectangle}, and drag a rectangle. Repeat \gui{Rectangular Zooming} for each further zoom. \\
    Freeze or release the view & \gui{Lock zooming and panning} disables wheel zoom and panning; \gui{Unlock zooming and panning} re-enables them. \\
    Cursor readout & Move the pointer over the plot; time, frequency and intensity appear in the status bar. \\
    Estimate drift rate & Click points along the burst; right-click or double-click to finish (Section~\ref{sec:drift}). \\
    Isolate burst & Press, drag a loop around the burst and release (Section~\ref{sec:isolation}). \\
    Ruler & Click two points (Section~\ref{sec:ruler}). \\
    Light curve by click & With click mode on, click a frequency (Section~\ref{sec:light-curves}). \\
    Annotations & Line: click the start and end. Text: click the position. Polygon: click the vertices and right-click to close (Section~\ref{sec:annotations}). \\
    \bottomrule
  \end{tabularx}
\end{table}

\gui{Reset Selection}\index{Reset Selection} clears an isolated-burst selection and drift markers;
\gui{Reset to Raw}\index{Reset to Raw} discards all processing; \gui{Reset All}\index{Reset All} clears the workspace
(Table~\ref{tab:toolbar}).

\section{Setting an exact display range}
\label{sec:display-range}
\index{display range}

\menu{View > Set Display Range...}\index{display range!dialog} sets the visible window numerically, which is useful
for comparing several stations over identical axes (Figure~\ref{fig:display-range}).
\begin{itemize}
  \item Choose \gui{Use seconds from file start} or \gui{Use UT clock time}, then enter the
    start and stop times in the \gui{Seconds} or \gui{UT} group. UT input is unavailable when
    the file has no \texttt{TIME-OBS}\index{TIME-OBS keyword} value.
  \item Enter the lower and upper limits in the \gui{Frequency} group (MHz).
  \item The full time and frequency ranges of the data are shown for reference.
\end{itemize}

\screenshot[0.5\linewidth]{dialogs/display_range_dialog}{The \gui{Set Display Range} dialog.}{fig:display-range}{The Set Display Range dialog with the Seconds/UT choice, the UT time fields and the Frequency fields, and the full-range information lines.}

\subsection{Display-range presets}
\index{display range!presets}

A display range can be stored under a name with \menu{View > Save Display Range
Preset...}, applied later to another file with \menu{Apply Display Range Preset\index{display range!presets}...}, and
removed with \menu{Delete Display Range Preset...}. Presets are stored in your preferences
and are available in every session.

\subsection{View configurations}
\label{sec:view-config}
\index{view configuration}\index{efaview.json@.efaview.json file}

\menu{View > Export View Config...} saves the display range, units, thresholds, color map
and graph styling to a portable \file{.efaview.json}\index{efaview.json@.efaview.json file} file; \menu{View > Import View
Config...} applies such a file. View configurations can be shared with colleagues and are
also accepted by the batch processor (Section~\ref{sec:batch}) and the Multi-Station
Comparison (Chapter~\ref{ch:comparison}), so that many spectra can be rendered with
identical axes and appearance.

\section{RFI cleaning}
\label{sec:rfi}
\index{RFI cleaning}

Radio-frequency interference appears as bright horizontal streaks (persistent
narrow-band emitters) and short spikes. \menu{Processing > RFI Cleaning} offers
\gui{Open RFI Panel}, \gui{Apply RFI} and \gui{Reset RFI}. The \gui{RFI Cleaning} panel
(Figure~\ref{fig:rfi}) applies a deterministic pipeline to the dynamic spectrum:
\begin{steps}
  \item two-dimensional median smoothing with a window of
    \gui{Kernel (freq)}~\(\times\)~\gui{Kernel (time)}\index{RFI cleaning!kernel};
  \item detection of hot channels from a robust \(Z\)-score of each channel's level and
    variability;
  \item repair of the masked channels from their neighbors;
  \item per-channel clipping of high outliers above an upper percentile (the low side is
    preserved).
\end{steps}

\screenshot[0.5\linewidth]{dialogs/rfi_cleaning_dialog}{The \gui{RFI Cleaning} panel.}{fig:rfi}{The RFI Cleaning panel with Enable RFI cleaning, the kernel, Z-threshold and percentile fields, the Masked channels line after Preview, and the Preview/Apply/Reset/Close buttons.}

\begin{table}[!htbp]
  \centering
  \caption{RFI cleaning parameters.}
  \label{tab:rfi}
  \small
  \begin{tabularx}{\linewidth}{@{}P{0.21\linewidth}P{0.15\linewidth}L@{}}
    \toprule
    \thead{Parameter} & \thead{Range (default)} & \thead{Effect} \\
    \midrule
    Enable RFI cleaning & on/off (on) & Includes RFI cleaning in the processing chain. \\
    Kernel (time) & 1--31, odd (3) & Median window along time. Larger values suppress short spikes more strongly but can blur fast burst structure. \\
    Kernel (freq) & 1--31, odd (3) & Median window across frequency channels. Larger values reduce narrow-band striping but can widen spectral features. \\
    Channel Z threshold\index{RFI cleaning!channel Z threshold} & 0.5--20 (6.0) & Cut-off for hot-channel masking. Lower values mask more channels. \\
    Percentile clip\index{RFI cleaning!percentile clip} & 90--99.99 (99.5) & Upper cap applied to each channel. Lower values clip peaks more aggressively. \\
    \bottomrule
  \end{tabularx}
\end{table}

\begin{description}
  \item[Preview] shows the result without committing it. The \gui{Masked channels}\index{RFI cleaning!masked channels} line
    reports the number and indices of the channels flagged as hot.
  \item[Apply] commits the cleaning to the processed data.
  \item[Reset] restores the default parameters in the panel.
  \item[Close] closes the panel.
\end{description}

If channel streaks remain, lower the \gui{Channel Z threshold} gradually (for example from
6.0 to 5.0 and then 4.0). If burst detail looks over-smoothed, reduce the kernel sizes or
raise the threshold. \menu{Edit > Reset to Raw} reverts all applied processing, including
RFI cleaning.

\section{Processing presets}
\label{sec:presets}
\index{presets}

A processing preset stores the display thresholds and their scale, the background
subtraction method, the intensity and time units, the color map, the graph styling, the
RFI settings and the default annotation style. Presets are managed from
\menu{Processing > Presets}:
\begin{description}
  \item[Raw FITS Percentile (5-98\%)] sets the threshold limits to the 5th and 98th
    percentiles of the data currently shown --- a quick starting range for raw files.
  \item[Save Current as Preset...] stores the current settings under a name.
  \item[Apply Preset...] applies a stored preset to the loaded data.
  \item[Set Default Preset...] chooses a preset that is applied automatically to every
    file loaded afterwards.
  \item[Clear Default Preset] stops applying a default preset.
  \item[Delete Preset...] removes a stored preset.
\end{description}

\section{Annotations}
\label{sec:annotations}
\index{annotations}

Annotations mark features on the spectrum. They are drawn in both renderers, saved in
projects, and included in exported figures and reports. \menu{Processing > Annotations}
(Figure~\ref{fig:annotation-menu}) contains:
\begin{description}
  \item[Add Polygon] click the vertices; right-click to close the polygon.
  \item[Add Line] click the start point and then the end point.
  \item[Add Text] click where the label should go, then enter it in the text dialog.
  \item[Edit Text Label...] choose an existing label and change its text and style.
  \item[Move Text Label...] choose a label, then click its new position.
  \item[Toggle Visibility] show or hide all annotations.
  \item[Delete Last] remove the most recent annotation.
  \item[Clear All] remove every annotation.
\end{description}
The text dialog sets the \gui{Text}, \gui{Font}, \gui{Font Size}, \gui{Font Style}
(\gui{Bold}, \gui{Italic}) and \gui{Color} of a label.

\screenshot[0.24\linewidth]{dialogs/annotation_text_dialog}{The \menu{Processing > Annotations} menu.}{fig:annotation-menu}{The text annotation dialog with the Text, Font, Font Size, Font Style and Color rows.}

\section{Light curves}
\label{sec:light-curves}
\index{light curve}

A light curve is the intensity at one frequency as a function of time. Light curves are
drawn on top of the spectrum (Figure~\ref{fig:light-curves}) from
\menu{Analysis > Plot Light Curves}:
\begin{description}
  \item[Enter frequency] opens a dialog for typing a frequency in MHz; the available range
    of the data is shown.
  \item[Click on a frequency] toggles click mode. While it is on, clicking the spectrum
    plots the light curve at the clicked frequency.
  \item[Settings...] opens the \gui{Light Curve Settings}\index{light curve!settings} dialog (Table~\ref{tab:light-curves}).
  \item[Clear light curve(s)] removes all light curves without affecting the data.
\end{description}

\screenshot{dialogs/light_curve_overlay}{Light curves at two frequencies overlaid on a dynamic spectrum, with the \menu{Analysis > Plot Light Curves} submenu open.}{fig:light-curves}{A background-subtracted spectrum with two light curves at different frequencies, frequency labels shown, and the Light Curve Settings dialog open beside it.}

\begin{table}[!htbp]
  \centering
  \caption{Light-curve settings.}
  \label{tab:light-curves}
  \small
  \begin{tabularx}{\linewidth}{@{}P{0.2\linewidth}P{0.22\linewidth}L@{}}
    \toprule
    \thead{Setting} & \thead{Range} & \thead{Effect} \\
    \midrule
    Mode & Single / Multiple & \gui{Single light curve} replaces the previous curve; \gui{Multiple light curves} keeps adding curves. \\
    Color & color picker & Color of the curve. \\
    Thickness & 0.5--12 & Line width. \\
    Vertical Scale & 0.1--5 & Height of the curve relative to the plot. \\
    Opacity & 0.1--1 & Transparency of the curve. \\
    Line Style & solid, dashed, dotted & Line pattern. \\
    Labels & on/off & \gui{Show frequency labels} next to each curve. \\
    \bottomrule
  \end{tabularx}
\end{table}
\gui{Apply} applies the settings, \gui{Reset to Defaults} restores the defaults and
\gui{Close} closes the dialog. Light curves are saved in projects and can be included in
project reports.

\section{Estimating the drift rate}
\label{sec:drift}
\index{drift rate!estimate}

The \gui{Estimate Drift Rate}\index{drift rate!estimate tool} toolbar tool gives a quick estimate of a burst's frequency
drift from points that you click.
\begin{steps}
  \item Click \gui{Estimate Drift Rate}.
  \item Click two or more points along the burst.
  \item Right-click or double-click to finish.
\end{steps}
The points are joined by dashed segments, and the status bar shows the
\gui{Average Drift Rate} in MHz/s (the mean of the segment drift rates), the start and end
frequencies, and the duration. Segments of zero duration are ignored. \gui{Reset Selection}
clears the drift markers. For a fitted drift rate and shock parameters, use the Type~II
analysis tools (Chapter~\ref{ch:type-ii}).

\section{Ruler measurement}
\label{sec:ruler}

\menu{Analysis > Ruler Measurement} measures between two points clicked on the spectrum.
The readout appears in the sidebar's \gui{Ruler Measurement} panel
(Section~\ref{sec:ruler-panel}): the coordinates of both points, the duration, the
frequency change and the slope. \menu{Analysis > Clear Ruler Measurement} or the panel's
\gui{Clear} button removes it. The ruler is kept when the timeline is extended.

\section{Overlays from other instruments}

Two overlays add context directly to the main plot:
\begin{itemize}
  \item the GOES X-ray flux curves, enabled from \menu{Solar Events > GOES Overlay}\index{GOES!overlay}
    (Section~\ref{sec:goes-overlay}); and
  \item the STEREO/SWAVES spectrum below the CALLISTO spectrum
    (Section~\ref{sec:swaves}).
\end{itemize}
